% Bản nháp chương 1 theo dàn bài mới. Chưa nối vào report.tex.
% ===========================================================================
\chapter{Bài toán định giá và hàm mục tiêu}
\label{sec:problem}

Bài toán ứng dụng của báo cáo là định giá một chiếc điện thoại đã qua sử dụng từ
các đặc điểm quan sát được của nó: hãng máy, dung lượng RAM, tuổi máy, tình trạng
màn hình, giá lúc mới mua. Mô hình đơn giản nhất cho mỗi đặc điểm một trọng số rồi
cộng lại, nên việc định giá quy về việc tìm bộ trọng số $w$ sao cho dự đoán sát
giá thật nhất. Tìm $w$ chính là giải một bài toán tối ưu hóa, và đó là nội dung
của báo cáo này.

\section{Bài toán ứng dụng và cách đo một lời giải tốt}
\label{sec:application}

Bộ dữ liệu \emph{Used Phone Price Prediction} trên Kaggle gồm \num{1000000} bản
ghi và 28 cột, phân bố theo bảng~\ref{tab:dataset}, trong đó không cột nào có giá
trị thiếu và không bản ghi nào trùng lặp. Vì bản đăng trên Kaggle đã qua tiền xử
lý, phần chuẩn bị dữ liệu ở chương~\ref{sec:data} bỏ được cả bước điền khuyết lẫn
bước khử trùng lặp. Thuận lợi này gắn với nguồn dữ liệu chứ không gắn với bài
toán: nếu thu thập trực tiếp từ các tin rao bán thì hai con số 0 nói trên không
còn, và nhóm dự đoán khi đó chi phí làm sạch sẽ vượt hẳn chi phí của phần tối ưu
hóa.

\begin{table}[tp]
  \centering
  \caption{Thành phần 28 cột của bộ dữ liệu gốc.}
  \label{tab:dataset}
  \begin{tabular}{lrl}
    \toprule
    Nhóm cột & Số cột & Ví dụ \\
    \midrule
    Định lượng liên tục & 14 & dung lượng RAM, tuổi máy, giá gốc, điểm camera \\
    Định lượng nhị phân & 7  & \texttt{screen\_cracked}, \texttt{box\_available} \\
    Định tính           & 6  & \texttt{brand} 7 mức, \texttt{model} 23 mức \\
    Biến mục tiêu       & 1  & \texttt{resale\_price} \\
    \bottomrule
  \end{tabular}
\end{table}

Giá bán lại lệch phải rõ rệt, với trung vị \num{18555} thấp hơn hẳn trung bình
\num{22598} trong khi giá trị lớn nhất lên tới \num{160257}. Nhóm vì thế lấy
logarit biến mục tiêu trước khi hồi quy: giữ nguyên thang gốc thì bình phương sai
số để vài trăm bản ghi đắt tiền chi phối hàm mục tiêu, và nghiệm thu được sẽ phản
ánh riêng nhóm máy đó thay vì phản ánh cả thị trường.

Phép lấy logarit đổi luôn cách đọc sai số, nên cần nói rõ ngay từ đây vì mọi con
số về chất lượng dự đoán trong báo cáo đều đọc theo cách này. Sai số bình phương
trung bình trên thang logarit, viết tắt là RMSE, đo độ lệch nhân tính chứ không đo
độ lệch cộng tính. Một mô hình đạt RMSE $\num{0.20601}$ tương ứng với hệ số lệch
điển hình $e^{0.20601} = \num{1.2288}$, tức dự đoán vênh khoảng \num{22.9}\% so
với giá thật, và trên một chiếc máy giá trung vị \num{18555} thì con số đó là
khoảng \num{4245} đơn vị tiền. Cách quy đổi này chỉ gần đúng, vì RMSE lấy trung
bình của bình phương nên nó nặng hơn độ lệch trung bình, nhưng nó đủ để so hai mô
hình với nhau và đó là việc duy nhất báo cáo cần tới.

Toàn bộ thí nghiệm chạy trên một mẫu ngẫu nhiên \num{200000} bản ghi, chia thành
\num{160000} điểm huấn luyện và \num{40000} điểm kiểm tra. Nhóm lấy mẫu vì chi phí
tính toán chứ không vì lý do thống kê: trên toàn bộ một triệu bản ghi, một lần
tính gradient mất \SI{271}{\milli\second} thay vì \SI{39}{\milli\second}, và vì
lưới tham số ở chương~\ref{sec:mechanisms} nhân con số đó lên theo tổng số vòng
lặp của mọi cấu hình nên chênh lệch gần bảy lần ấy nhân lên theo toàn bộ lưới. Cỡ
mẫu này đủ cho các kết luận về tốc độ hội tụ, vốn chỉ phụ thuộc $L$, $\mu$ và
$\kappa$, nhưng sẽ không đủ nếu đại lượng cần đo phụ thuộc trực tiếp vào $n$,
chẳng hạn phương sai gradient của SGD ở chương~\ref{sec:mechanisms}.

\section{Từ định giá tới hàm mục tiêu Ridge}
\label{sec:objective}

Ký hiệu $n$ là số điểm dữ liệu, $d$ là số thuộc tính sau khi mã hóa, $X \in
\R^{n \times d}$ là ma trận thiết kế và $y \in \R^{n}$ là vector logarit giá. Sau
khi chuẩn hóa từng cột của $X$ và trừ trung bình khỏi $y$ thì hệ số chặn tối ưu
bằng không, nên bài toán chỉ còn một khối biến $w$:
%
\begin{equation}
  \label{eq:objective}
  f(w) = \frac{1}{2n} \norm{Xw - y}^{2} + \frac{\lambda}{2} \norm{w}^{2}.
\end{equation}
%
Công thức này ghi lại một đánh đổi. Số hạng thứ nhất ép mô hình bám dữ liệu, số
hạng thứ hai ép các trọng số nhỏ lại, và $\lambda$ định giá quy đổi giữa hai đòi
hỏi đó. Hai hằng số phía trước có mặt vì lý do kỹ thuật chứ không vì lý do thống
kê: hệ số $\tfrac{1}{2}$ để đạo hàm không kéo theo số 2, còn hệ số $\tfrac{1}{n}$
để hằng số Lipschitz ở \eqref{eq:constants} không phụ thuộc cỡ mẫu, nhờ vậy mọi
kết luận về độ dài bước vẫn dùng được khi đổi số điểm dữ liệu. Bỏ $\tfrac{1}{n}$
đi thì $L$ tăng theo $n$ và mỗi lần đổi cỡ mẫu lại phải chọn lại độ dài bước.

Đạo hàm của \eqref{eq:objective} tính trực tiếp được:
%
\begin{equation}
  \label{eq:derivatives}
  \nabla f(w) = \frac{1}{n} X\trans (Xw - y) + \lambda w,
  \qquad
  \nabla^{2} f(w) = \frac{1}{n} X\trans X + \lambda I.
\end{equation}
%
Điểm đáng chú ý nằm ở vế thứ hai: Hessian không chứa $w$, tức độ cong của $f$
giống nhau tại mọi điểm. Tính chất này là lý do phần lớn báo cáo đơn giản đi so
với một bài toán phi tuyến, và chương~\ref{sec:algorithms} sẽ cho thấy nó kéo theo
việc phương pháp Newton hội tụ chỉ sau một vòng lặp.

Vì $\lambda > 0$ nên Hessian xác định dương và khả nghịch, do đó $f$ lồi mạnh và
có đúng một điểm cực tiểu, giải ra được bằng cách cho gradient bằng không:
%
\begin{equation}
  \label{eq:closed-form}
  \wstar = \left( \frac{1}{n} X\trans X + \lambda I \right)^{-1}
           \left( \frac{1}{n} X\trans y \right),
  \qquad \fstar = f(\wstar).
\end{equation}
%
Hệ quả của \eqref{eq:closed-form} quan trọng hơn bản thân công thức. Vì $\fstar$
tính được trước khi chạy bất kỳ thuật toán nào nên sai số $f(w_k) - \fstar$ đo
được tuyệt đối, không phải so với giá trị tốt nhất quan sát được như ở các bài
toán không có nghiệm đóng. Mọi biểu đồ hội tụ trong báo cáo dựa vào đại lượng này,
và hệ số co rút thực nghiệm ở chương~\ref{sec:theory} đối chiếu thẳng được với cận
lý thuyết cũng nhờ nó.

\section{Ba con số quyết định độ khó}
\label{sec:constants}

Tốc độ của mọi phương pháp bậc một trong báo cáo phụ thuộc vào ba đại lượng:
%
\begin{equation}
  \label{eq:constants}
  L = \lambda_{\max}\!\Bigl( \tfrac{1}{n} X\trans X \Bigr) + \lambda,
  \qquad
  \mu = \lambda_{\min}\!\Bigl( \tfrac{1}{n} X\trans X \Bigr) + \lambda,
  \qquad
  \kappa = \frac{L}{\mu}.
\end{equation}
%
$L$ là hằng số Lipschitz của gradient và $\mu$ là hệ số lồi mạnh. Với một hàm bậc
hai, hai đại lượng ấy đọc thẳng từ phổ của Hessian: gradient thỏa
$\norm{\nabla f(u) - \nabla f(v)} \le L \norm{u - v}$ với $L$ chính là chuẩn phổ
của $\nabla^{2} f$, vì Hessian hằng theo \eqref{eq:derivatives} nên chặn này đạt
được và không thể lấy nhỏ hơn.

Số điều kiện $\kappa$ có một cách đọc hình học giải thích gần hết kết quả về sau.
Các đường mức của $f$ là những ellipsoid, và $\kappa$ đúng bằng tỉ số giữa trục
dài nhất với trục ngắn nhất của chúng. Khi $\kappa = 1$, đường mức là mặt cầu và
gradient chỉ thẳng về tâm, nên một bước duy nhất là tới đích. Khi $\kappa$ lớn,
ellipsoid bị bóp dẹt thành một cái máng hẹp và dài, và gradient khi đó gần như
vuông góc với trục của máng. Sai số theo chiều ngang máng vì thế giảm rất nhanh,
còn sai số dọc theo máng giảm chậm hơn nhiều bậc, và chính thành phần thứ hai
quyết định số vòng lặp. Các cơ chế ở chương~\ref{sec:mechanisms} đều quy về hình
ảnh này.

\begin{table}[tp]
  \centering
  \caption{Các hằng số của bài toán sau khi cố định dữ liệu và $\lambda$.}
  \label{tab:constants}
  \begin{tabular}{llr}
    \toprule
    Ký hiệu & Ý nghĩa & Giá trị \\
    \midrule
    $n$       & số điểm huấn luyện              & \num{160000} \\
    $d$       & số thuộc tính sau khi mã hóa    & 280 \\
    $\lambda$ & hệ số hiệu chỉnh                & \num{0.01} \\
    $L$       & hằng số Lipschitz của gradient  & \num{2.68309} \\
    $\mu$     & hệ số lồi mạnh                  & \num{0.01} \\
    $\kappa$  & số điều kiện                    & \num{268.309} \\
    $\fstar$  & giá trị tối ưu                  & \num{0.0237430} \\
    \bottomrule
  \end{tabular}
\end{table}

Bảng~\ref{tab:constants} cho thấy $\mu$ trùng khít với $\lambda$, và đó không phải
trùng hợp. Ma trận Gram $\tfrac{1}{n} X\trans X$ có hạng 273 trên 280, tức bảy
hướng hoàn toàn suy biến, và ngay cả trị riêng đầu tiên sau bảy hướng đó cũng chỉ
cỡ $2 \cdot 10^{-8}$, nhỏ hơn $L$ tới tám bậc. Dữ liệu tự nó có những hướng phẳng
lì, và thứ duy nhất tạo ra độ cong theo các hướng ấy là số hạng phạt, nên
%
\begin{equation}
  \label{eq:mu-equals-lambda}
  \mu = \lambda, \qquad \kappa = \frac{L}{\lambda}.
\end{equation}
%
Độ khó của bài toán do đó là thứ người làm chọn qua $\lambda$, không phải thứ dữ
liệu áp đặt, và chương~\ref{sec:lambda} khai thác đúng tính chất này. Quan hệ
\eqref{eq:mu-equals-lambda} đảo chiều nếu mọi cột độc lập tuyến tính và trị riêng
nhỏ nhất cùng bậc với $\lambda$; khi đó $\mu$ quay về do dữ liệu quyết định và
$\lambda$ mất vai trò điều khiển $\kappa$.

\section{Năm thuật toán của báo cáo}
\label{sec:catalogue}

Bảng~\ref{tab:methods} giới thiệu năm thuật toán mà chương~\ref{sec:compare} đem
so sánh, đủ để đọc chương đó trước khi tới phần cài đặt chi tiết ở
chương~\ref{sec:algorithms}. Hai cột cuối là hai thứ phân biệt chúng: mỗi phương
pháp dùng bao nhiêu thông tin về $f$, và trả giá bao nhiêu cho mỗi vòng lặp.

\begin{table}[tp]
  \centering
  \caption{Năm thuật toán trong mạch chính, xếp theo lượng thông tin sử dụng.}
  \label{tab:methods}
  \begin{tabular}{llll}
    \toprule
    Thuật toán & Thông tin dùng & Chi phí mỗi vòng & Tham số phải chọn \\
    \midrule
    Gradient descent & gradient tại $w_k$      & $\bigO(nd)$        & $t$, hoặc $\alpha$ và $\beta$ \\
    SGD              & gradient trên lô $B$ điểm & $\bigO(Bd)$      & $B$, quy tắc chọn $\eta_k$ \\
    AGD (Nesterov)   & gradient tại điểm ngoại suy & $\bigO(nd)$    & $t$, công thức momentum \\
    L-BFGS           & $m$ cặp hiệu gần nhất   & $\bigO(nd + md)$   & $m$ \\
    Newton           & toàn bộ Hessian         & $\bigO(nd^{2} + d^{3})$ & không có \\
    \bottomrule
  \end{tabular}
\end{table}

Cột chi phí giải thích trước một điều mà chương~\ref{sec:compare} sẽ đo: xếp hạng
theo số vòng lặp và xếp hạng theo thời gian chạy là hai thứ khác nhau. Newton dùng
nhiều thông tin nhất nên cần ít vòng lặp nhất, nhưng mỗi vòng của nó đắt hơn một
vòng gradient descent theo bậc $d^{2}$, và với $d = 280$ thì chênh lệch ấy đủ lớn
để đảo thứ hạng theo hướng nào cũng có thể xảy ra. Nhóm còn cài Heavy ball,
Newton-CG và Adam nhưng không đưa vào so sánh chính; lý do nằm ở phụ lục~\ref{app:cut}.

\section{Câu hỏi báo cáo trả lời}
\label{sec:question}

Bài toán này có nghiệm đóng ở \eqref{eq:closed-form}, tính được bằng một lần phân
rã Cholesky. Câu hỏi tự nhiên là vậy chạy các thuật toán lặp để làm gì.

Có hai câu trả lời, và chúng định hình toàn bộ phần còn lại. Thứ nhất, chính vì có
nghiệm đóng mà bài toán này dùng được làm phòng thí nghiệm: $\fstar$ biết trước
tới sai số máy, nên mọi phát biểu về tốc độ hội tụ đo được bằng một cái thước
tuyệt đối, thứ không tồn tại ở các bài toán mà người ta thực sự cần phương pháp
lặp. Thứ hai, nghiệm đóng chỉ dùng được trong một khoảng quy mô nhất định: chi phí
$\bigO(d^{3})$ và việc phải lập ma trận $d \times d$ trở thành rào cản khi $d$ lớn,
và khi đó những phương pháp mà báo cáo khảo sát mới là lựa chọn duy nhất.

Câu hỏi xuyên suốt báo cáo vì thế phát biểu được thành một dòng: \emph{huấn luyện
mô hình định giá là giải một bài toán tối ưu hóa; giải bằng cách nào, và giải tới
mức nào thì dừng?} Bảng~\ref{tab:roadmap} chỉ ra chương nào trả lời phần nào của
câu hỏi đó.

\begin{table}[tp]
  \centering
  \caption{Câu hỏi mà mỗi chương trả lời.}
  \label{tab:roadmap}
  \begin{tabular}{ll}
    \toprule
    Chương & Câu hỏi \\
    \midrule
    \ref{sec:data}       & Chuẩn bị dữ liệu làm bài toán dễ hay khó đi, đo bằng $\kappa$ \\
    \ref{sec:enough}     & Sai số $f - \fstar$ nhỏ tới đâu thì sai số giá ngừng giảm \\
    \ref{sec:compare}    & Thuật toán nào về đích trước, kể cả khi đặt cạnh thư viện \\
    \ref{sec:algorithms} & Năm thuật toán được cài đặt và đo thời gian ra sao \\
    \ref{sec:mechanisms} & Ba cơ chế giải thích vì sao các đường nằm ở đó \\
    \ref{sec:theory}     & Cận lý thuyết mô tả đúng chỗ nào và lệch chỗ nào \\
    \ref{sec:lambda}     & Đổi $\lambda$, tức đổi $\kappa$, thì thứ hạng có đổi không \\
    \ref{sec:conclusion} & Rút lại thì nên chọn thuật toán nào, với điều kiện gì \\
    \bottomrule
  \end{tabular}
\end{table}

Kết quả chính, nêu trước ở đây và trình bày đầy đủ ở chương~\ref{sec:compare}: năm
cấu hình tốt nhất của năm thuật toán cho RMSE trên tập kiểm tra trùng nhau tới chữ
số thứ bảy, trong khi thời gian để cùng đạt mức sai số $10^{-6}$ trải hơn một bậc,
từ \SI{0.076}{\second} của Newton tới \SI{0.879}{\second} của AGD. Chọn sai thuật
toán vì thế không làm cửa hàng định giá sai, nó chỉ làm người huấn luyện chờ lâu
hơn. Ngoại lệ duy nhất là SGD với độ dài bước hằng, dừng ở RMSE $\num{0.20986}$ so
với $\num{0.20601}$, tức đắt thêm khoảng 88 đơn vị tiền trên mỗi chiếc máy giá
trung vị. Chương~\ref{sec:enough} truy nguyên nhân của ngoại lệ ấy, và
chương~\ref{sec:compare} trình bày đầy đủ phần so sánh.
